\section{Related Work}
% deception detection - text etc 
% nlp + deception - datasets mainly specifically nlp venue
% models tried by pople

\paragraph{Deception detection.} Deception is an act emerging since the beginning of time with the Serpent and Eve in the Garden of Eden: \emph{And the serpent said unto the woman, Ye shall not surely die}. However, humans are often not very good at spotting these lies \cite{belot2012can,gneezy2005deception}, and are no better than making random decisions \cite{ockenfels2000experiment}. This raises the question of whether there are valid indicators of deception. According to studies \cite{wang2010pinocchio,depaulo2003cues,ekman1997deception,zuckerman1981verbal, wang2010pinocchio}, employing non-verbal signals such as visual, facial, and aural cues might significantly help distinguish sincere and opportunistic communication. 
Recent studies mention participants are significantly more accurate at spotting lies from both audio and videos (82\%) or only videos (66\%), compared to text (57\%) \cite{wittenberg2021minimal,groh2022human}. In this paper, we investigate if a computational model can instead detect deception in text using language cues.


\paragraph{Deception + NLP datasets.} Automated deception detection techniques so far predominantly utilized visual cues such as facial or eye movements to detect deception and time to response (\cite{meservy2005deception,gonzalez2019can}, or linguistic cues from transcriptions from court hearings \cite{fornaciari2013automatic}, deceptive hotel reviews \cite{ott2013negative}, news articles from Buzzfeed datasets \cite{potthast2017stylometric}, and fact-checked tweets \cite{van2022personal}. Datasets that require assigning specific individuals roles (lier/truth-teller) include a multimodal conversational dataset Box of Lies \cite{soldner2019box}, Golden Globes \cite{darai2013attraction} differ from Diplomacy gameplay \cite{peskov2020takes}, and Real or Spiel \cite{ho2016computer} where one can \emph{choose} to lie. In contrast, our derived dataset from this game show is conversational, grounded in a real deceptive environment with the presence of objective truth to detect the deception not present in existing datasets. The only work, \cite{banerjee2023experience} that investigated, To Tell The Truth gameshow neither built computational models nor analyzed cues from interactions. 

\paragraph{Models detecting deception.}
Computational models that focus on language cues mainly use psycholinguistic features \cite{Grlea2016PsycholinguisticFF, soldner2019box}, or syntactic parse of the texts \cite{soldner2019box} to identify predictors for deception. \citet{ho2016computer} utilized a power dynamics vocabulary to identify deception in long-term relationships. In our paper, we focus on more complex signals, such as ambiguity or half-truths, and benchmark the performance of LLMs for the first time in the context of deception detection from text.